% GENERATED FILE. Edit canonical lessons, then run npm run generate:books.
% canonical-source-hash: fnv1a64:98c005b09b835e92
% canonical-lessons: SA-C450-visadah, SA-C450-socati, SA-C450-rodanam, SA-C450-baspah, SA-C450-nihsvasah

\chapter{Sadness and Tears}
\label{ch:sa-sadness-and-tears}
\hlchaptermodality{450}

\begin{chapteropening}
\textbf{By the end of this chapter:} \emph{I can say that someone was sad and why.}

Complete the last lesson of chapter 450: I can say that someone was sad and why.
\end{chapteropening}

\section[\texorpdfstring{\sk{विषादः} (viṣādaḥ)}{विषादः (viṣādaḥ)}]{\sk{विषादः} (viṣādaḥ) — dejection}

\label{lesson:SA-C450-visadah}



\begin{quote}
Before the new one: say the Sanskrit for pleases, then the Sanskrit for a thrill.
\end{quote}

\subsection*{You'll want to know: \sk{विषादः}}
\textbf{\sk{विषादः}} — \emph{viṣādaḥ} — \textquotedblleft{}dejection\textquotedblright{}.

\begin{quote}
\sk{परीक्षायाः} \sk{परिणामेन} \sk{छात्रस्य} \sk{विषादः} \sk{अभवत्।}
\end{quote}

\begin{quote}
\emph{parīkṣāyāḥ pariṇāmena chātrasya viṣādaḥ abhavat.}
\end{quote}

\begin{quote}
The student was dejected by the exam result.
\end{quote}

A masculine noun.

\begin{cousinweb}
\textbf{\sk{विषादः}} (\emph{viṣādaḥ}) is built from \emph{vi}, \textquotedblleft{}apart\textquotedblright{}, and the root √sad, \textquotedblleft{}to sit, to sink\textquotedblright{}: a sinking of the heart; MW gives \textquotedblleft{}dejection, depression, despondency\textquotedblright{} (MW).
\end{cousinweb}

\begin{grammarlens}[title={What you've built}]
One more word for describing what you lived and felt.
\end{grammarlens}

\subsection*{Your turn}
\begin{itemize}
\raggedright
  \item \emph{Say it:} \emph{viṣādaḥ}
  \item \emph{Say it:} \emph{viṣādaḥ}, once more
  \item \emph{Say it:} \emph{parīkṣāyāḥ pariṇāmena chātrasya viṣādaḥ abhavat.} — The student was dejected by the exam result.
  \item \emph{Recall:} say the Sanskrit for pleases, then the Sanskrit for a thrill, then say \emph{viṣādaḥ} again
\end{itemize}

\begin{tcolorbox}[breakable,colback=teal!4,colframe=teal!35!black,title={Before you move on}]
What does \sk{विषादः} mean? (\textquotedblleft{}Dejection\textquotedblright{}.) Say it once more.
\end{tcolorbox}
\section[\texorpdfstring{\sk{शोचति} (śocati)}{शोचति (śocati)}]{\sk{शोचति} (śocati) — grieves}

\label{lesson:SA-C450-socati}



\begin{quote}
Before the new one: say the Sanskrit for a thrill, then the Sanskrit for dejection.
\end{quote}

\subsection*{You'll want to know: \sk{शोचति}}
\textbf{\sk{शोचति}} — \emph{śocati} — \textquotedblleft{}grieves\textquotedblright{}.

\begin{quote}
\sk{बालिका} \sk{नष्टं} \sk{शुनकं} \sk{शोचति।}
\end{quote}

\begin{quote}
\emph{bālikā naṣṭaṃ śunakaṃ śocati.}
\end{quote}

\begin{quote}
The girl grieves for her lost dog.
\end{quote}

Past (he …): \sk{सः} \sk{बहु} \sk{शोचितवान्}, \emph{saḥ bahu śocitavān}.

\begin{cousinweb}
\textbf{\sk{शोचति}} (\emph{śocati}) comes from the root √śuc, \textquotedblleft{}to glow, to burn\textquotedblright{} and \textquotedblleft{}to grieve\textquotedblright{}, the root of \textbf{\sk{शोकः}} (MW).
\end{cousinweb}

\begin{grammarlens}[title={What you've built}]
One more word for describing what you lived and felt.
\end{grammarlens}

\subsection*{Your turn}
\begin{itemize}
\raggedright
  \item \emph{Say it:} \emph{śocati}
  \item \emph{Say it:} \emph{śocati}, once more
  \item \emph{Say it:} \emph{bālikā naṣṭaṃ śunakaṃ śocati.} — The girl grieves for her lost dog.
  \item \emph{Recall:} say the Sanskrit for a thrill, then the Sanskrit for dejection, then say \emph{śocati} again
\end{itemize}

\begin{tcolorbox}[breakable,colback=teal!4,colframe=teal!35!black,title={Before you move on}]
What does \sk{शोचति} mean? (\textquotedblleft{}Grieves\textquotedblright{}.) Say it once more.
\end{tcolorbox}
\section[\texorpdfstring{\sk{रोदनम्} (rodanam)}{रोदनम् (rodanam)}]{\sk{रोदनम्} (rodanam) — weeping; crying}

\label{lesson:SA-C450-rodanam}



\begin{quote}
Before the new one: say the Sanskrit for dejection, then the Sanskrit for grieves.
\end{quote}

\subsection*{You'll want to know: \sk{रोदनम्}}
\textbf{\sk{रोदनम्}} — \emph{rodanam} — \textquotedblleft{}weeping; crying\textquotedblright{}.

\begin{quote}
\sk{शिशोः} \sk{रोदनेन} \sk{माता} \sk{रात्रौ} \sk{जागर्ति।}
\end{quote}

\begin{quote}
\emph{śiśoḥ rodanena mātā rātrau jāgarti.}
\end{quote}

\begin{quote}
The baby's crying keeps its mother awake at night.
\end{quote}

A neuter noun.

\begin{cousinweb}
\textbf{\sk{रोदनम्}} (\emph{rodanam}) comes from the root √rud, \textquotedblleft{}to weep\textquotedblright{}, the root of \textbf{\sk{रोदिति}}; MW gives \textquotedblleft{}weeping, crying\textquotedblright{} and also \textquotedblleft{}tears\textquotedblright{} (MW).
\end{cousinweb}

\begin{grammarlens}[title={What you've built}]
One more word for describing what you lived and felt.
\end{grammarlens}

\subsection*{Your turn}
\begin{itemize}
\raggedright
  \item \emph{Say it:} \emph{rodanam}
  \item \emph{Say it:} \emph{rodanam}, once more
  \item \emph{Say it:} \emph{śiśoḥ rodanena mātā rātrau jāgarti.} — The baby's crying keeps its mother awake at night.
  \item \emph{Recall:} say the Sanskrit for dejection, then the Sanskrit for grieves, then say \emph{rodanam} again
\end{itemize}

\begin{tcolorbox}[breakable,colback=teal!4,colframe=teal!35!black,title={Before you move on}]
What does \sk{रोदनम्} mean? (\textquotedblleft{}Weeping; crying\textquotedblright{}.) Say it once more.
\end{tcolorbox}
\section[\texorpdfstring{\sk{बाष्पः} (bāṣpaḥ)}{बाष्पः (bāṣpaḥ)}]{\sk{बाष्पः} (bāṣpaḥ) — tears}

\label{lesson:SA-C450-baspah}



\begin{quote}
Before the new one: say the Sanskrit for grieves, then the Sanskrit for weeping.
\end{quote}

\subsection*{You'll want to know: \sk{बाष्पः}}
\textbf{\sk{बाष्पः}} — \emph{bāṣpaḥ} — \textquotedblleft{}tears\textquotedblright{}.

\begin{quote}
\sk{बालिकायाः} \sk{मुखे} \sk{बाष्पः} \sk{अस्ति।}
\end{quote}

\begin{quote}
\emph{bālikāyāḥ mukhe bāṣpaḥ asti.}
\end{quote}

\begin{quote}
There are tears on the girl's face.
\end{quote}

A masculine noun.

\begin{cousinweb}
\textbf{\sk{बाष्पः}} (\emph{bāṣpaḥ}) MW records \emph{bāṣpa}, also written \emph{vāṣpa}, \textquotedblleft{}a tear, tears\textquotedblright{} and \textquotedblleft{}steam, vapour\textquotedblright{}; it lives on as Hindi \emph{bhāp}, \textquotedblleft{}steam\textquotedblright{} (CDIAL 9223); MW gives no derivation, so its origin is uncertain (MW).
\end{cousinweb}

\begin{grammarlens}[title={What you've built}]
One more word for describing what you lived and felt.
\end{grammarlens}

\subsection*{Your turn}
\begin{itemize}
\raggedright
  \item \emph{Say it:} \emph{bāṣpaḥ}
  \item \emph{Say it:} \emph{bāṣpaḥ}, once more
  \item \emph{Say it:} \emph{bālikāyāḥ mukhe bāṣpaḥ asti.} — There are tears on the girl's face.
  \item \emph{Recall:} say the Sanskrit for grieves, then the Sanskrit for weeping, then say \emph{bāṣpaḥ} again
\end{itemize}

\begin{tcolorbox}[breakable,colback=teal!4,colframe=teal!35!black,title={Before you move on}]
What does \sk{बाष्पः} mean? (\textquotedblleft{}Tears\textquotedblright{}.) Say it once more.
\end{tcolorbox}
\section[\texorpdfstring{\sk{निःश्वासः} (niḥśvāsaḥ)}{निःश्वासः (niḥśvāsaḥ)}]{\sk{निःश्वासः} (niḥśvāsaḥ) — a sigh}

\label{lesson:SA-C450-nihsvasah}



\begin{quote}
Before the new one: say the Sanskrit for weeping, then the Sanskrit for tears.
\end{quote}

\subsection*{You'll want to know: \sk{निःश्वासः}}
\textbf{\sk{निःश्वासः}} — \emph{niḥśvāsaḥ} — \textquotedblleft{}a sigh\textquotedblright{}.

\begin{quote}
\sk{पिता} \sk{श्रान्तः}, \sk{सः} \sk{निःश्वासं} \sk{मुञ्चति।}
\end{quote}

\begin{quote}
\emph{pitā śrāntaḥ, saḥ niḥśvāsaṃ muñcati.}
\end{quote}

\begin{quote}
Father is tired; he lets out a sigh.
\end{quote}

A masculine noun.

\begin{cousinweb}
\textbf{\sk{निःश्वासः}} (\emph{niḥśvāsaḥ}) is built from \emph{niḥ}, \textquotedblleft{}out\textquotedblright{}, and the root √śvas, \textquotedblleft{}to breathe\textquotedblright{}, the root of \textbf{\sk{श्वासः}}: a breath let out (MW).
\end{cousinweb}

\begin{grammarlens}[title={What you've built}]
One more word for describing what you lived and felt.
\end{grammarlens}

\subsection*{Your turn}
\begin{itemize}
\raggedright
  \item \emph{Say it:} \emph{niḥśvāsaḥ}
  \item \emph{Say it:} \emph{niḥśvāsaḥ}, once more
  \item \emph{Say it:} \emph{pitā śrāntaḥ, saḥ niḥśvāsaṃ muñcati.} — Father is tired; he lets out a sigh.
  \item \emph{Recall:} say the Sanskrit for weeping, then the Sanskrit for tears, then say \emph{niḥśvāsaḥ} again
\end{itemize}

\begin{tcolorbox}[breakable,colback=teal!4,colframe=teal!35!black,title={Before you move on}]
What does \sk{निःश्वासः} mean? (\textquotedblleft{}A sigh\textquotedblright{}.) Say it once more.
\end{tcolorbox}
